
\section{ベイズの定理}
ここまで尤度を使った推定方法について説明してきました。モーメント法と最尤法，いずれも同じ答えが出るのですがその考え方が違うというところを確認してきましたね。ここでは最後に，もう1つの推定法であるベイズ法について解説します。そのためには，ベイズ法の根本的な法則について知っておいた方が良いので，少し回りくどいかもしれませんが，まずはそこから説明していきましょう\footnote{目的を先に教えてくれないとやだよ，という人のために一言で目的だけ言っておきます。ベイズ法は最尤法と同じく確率モデルを使うのですが，尤度が「確率じゃない」という妙な性質があったのに対し，ベイズ法は答えを確率で出します。尤度を確率にするための仕掛けがこれから始まるのです。}。
\subsection{ベイズの定理とは}
ベイズというのは人の名前で，トーマス・ベイズ（Thomas Bayes、1702年 - 1761年4月17日）はイギリスの長老派の牧師であり数学者でもある人です。この人の考えた「逆確率の法則」から端を発した推定法が，21世紀の現代においてとても重要なものになってきています。ベイズ牧師の考えついた法則は非常に単純なもので，簡単に書くと次のような式で表現できます。
\[ P(A|B) = \frac{P(B|A)P(A)}{P(B)} \]


例えば，ある学校の生徒数は男子が40人，女子が60人だったとします。そこからランダムに一人を選び出した時，それが男子である確率はどれぐらいでしょうか？ この確率を$P(\mbox{男子})$とすると，
\[P(\mbox{男子})=\frac{40}{(100)}=0.4 \]
であることはすぐにわかりますね。

では次に\keyterm{同時確率}{どうじかくりつ}{joint probability}という考え方を導入しましょう。これは先ほどの学校について，表\ref{tbl::24_02}のように学年情報を追加した例を使って説明します。


\begin{table}[ht]
    \centering
    \caption{ある学校の学年と性別のデータ} 
    \begin{tabular}{rrrrr}
      \hline
      & 1年 & 2年 & 3年 & 合計 \\
    \hline
    男子 & 15 & 12 & 13 & 40\\
    女子 & 22 & 20 & 18 & 60\\\hline
    合計 & 37 & 32 & 31 & 100\\\hline
    \end{tabular}
      \label{tbl::24_02}
\end{table}

ここでランダムに一人を選び出した時，それが「3年男子」である確率はどれぐらいでしょうか？ これは「学年は3」かつ「性別は男」なので，学年と性別の両方の条件が入った同時確率といい，$P(\mbox{学年},\mbox{性別})$と書きます。これは表から，次のように計算できますね。
\[ P(\mbox{学年;3},\mbox{性別;男})= \frac{13}{100}=0.13 \]

次に用語として\keyterm{周辺確率}{しゅうへんかくりつ}{marginal probability}というのを考えます。これはある分割に関して全て足し上げたもの，というものですが，要するに表\ref{tbl::24_02}の合計欄のことです。例えば「男子生徒についての周辺確率」は，男子生徒が学年によって3つに分割されていますが，これを全て足し上げるので，1年男子+2年男子+3年男子で求まるものです。なんてことはないですね。

では次のステップ。今度は\keyterm{条件つき確率}{じょうけんつきかくりつ}{conditional probability}というのを導入します。無作為に学生を選び出す時，「選ばれたのが女子である」ということが先にわかったとします\footnote{例えば名前とか容易に類推できた，といったようなシーンです。}。その時，その生徒が2年生である確率はどれくらいでしょうか？ この問いを「女子という条件がついた上での確率」という意味で，条件つき確率といい，表現としては$P(\mbox{2年生}|\mbox{女子})$と書きます。2つの分割を縦棒で区切り，条件の方が後ろに書きます。「女子という条件が与えられた上での二年生の確率は」というように読んだりします。

ちょっと状況はややこしくなりましたが，確率の計算そのものはそれほど難しくありませんね。女子である，というのが決まっているので，分母が変わるだけなのです。つまり，
\[P(\mbox{2年生}|\mbox{女子})=\frac{20}{60}=0.333... \]
となります。

条件つき確率の計算は分母が変わったと思えばいいのですね。この条件つき確率は，同時確率を周辺確率で割っても計算できます。すなわち，
\[ P(\mbox{2年生}|\mbox{女子}) = \frac{P(\mbox{2年生},\mbox{女子})}{P(\mbox{女子})}=\frac{0.2}{0.6}=0.333... \]

さて，これで用語の準備は終了です。ここから式の展開が面白くなってきます。記号で表現するために，性別の変数をAとし，$A_1=$男子，$A_2=$女子，学年の変数をBとして$B_1=$1年生，$B_2=$2年生，$B_3=$3年生としましょう。

あらためて，条件付き確率は同時確率を周辺確率で割ったもの，でしたね。つまり，
\[ P(B_j | A_i) = \frac{P(B_j,A_i)}{P(A_i)} \]
ここで右辺の分母を両辺にかけて，
\[ P(A_i)P(B_j|A_i) = P(B_j,A_i) \]
とします。これはA,Bが何であるかにかかわらず，一般的に成立します。たとえば$P(A_i|B_j)$と記号をひっくり返しても，
\[ P(A_i|B_j) = \frac{P(A_i,B_j)}{P(B_j)} \]
より
\[P(B_j)P(A_i|B_j) = P(A_i,B_j) \]
としても構いませんね。


ところで，「2年の男子」と「男子の2年」というのは同じ意味です。つまり，
\[P(B_j,A_i) = P(A_i, B_j) \]
ですから，上の式をつなぎ合わせて
\[P(A_i)P(B_j|A_i) = P(B_j)P(A_i|B_j) \]
とできます。ここで両辺を$P(B_j)$で割って清書すると，
\[ P(A_i | B_j )=\frac{P(B_j|A_i)P(A_i)}{P(B_j)} \]
が得られます。やったね！ベイズの定理が導出できました。


さて，これの何がすごいのでしょう？ 
\subsection{ベイズの定理の意味}
実は，\underline{左辺と右辺の中にある条件つき確率の中身がひっくり返っている}のがすごいんです。

凄さがピンとこないかもしれませんね。2年生だとか女子だとかだと，面白味がわからないかもしれませんので，別の言葉で考えてみましょうか。

コロナウイルスに感染しているかどうかは，PCR検査で明らかになりますが，これらの検査の結果というのも確率的な要素があるのはご存知でしょうか。本当は感染していないのに，検査では陽性になってしまう，ということがゼロではありません。こういうのを偽陽性といいます。逆に感染しているのに，検査で検出できずに陰性になる，偽陰性ということもあります。いずれも確率的な話です。

コロナに感染している$\to$検査で陽性になる，というのが因果関係ですが，我々は実際には検査で陽性である$\to$コロナに感染しているかもしれない，というように推論します。これを今日学んだ確率の言葉で書いてみましょう。

感染している状態で，陽性反応が出るというのは条件つき確率として表現できますね。$P(\mbox{陽性}|\mbox{感染})$，というわけです。さてここに，そもそもの感染率というのが別途見積もれたとします。$P(\mbox{感染})$ですね。さらにもうひとつ。感染しているかどうかに関係なく，陽性であるという結果だけ取り上げて$P(\mbox{陽性})$を考えることもできます。これが揃うと，先ほどの定理を使って
\[ \frac{P(\mbox{陽性}|\mbox{感染})P(\mbox{感染})}{P(\mbox{陽性})} = P(\mbox{感染}|\mbox{陽性})\]
を計算できます。つまり陽性になったという結果から，本当にかかっているかどうかの確率を計算できるのです。

多くの場合，検査で陽性だったら「もうだめだー」と思ってしまいがちです。これは陽性であると確実に検出できる，偽陽性がゼロの検査であれば正しい判断です。が，実際には確率的に判断すべき問題です。具体的に数字を例にして考えてみましょう。

\begin{itemize}
    \item 40歳の女性が乳がんにかかる確率は1\%です。
    \item 乳がんである人が乳房X線検査で陽性と出る確率は90\%です。
    \item 本当は乳がんでなくても，検査で陽性と出る確率は8\%にすぎません。
    \item あなたが40歳の女性で検査結果が陽性だったとしたら，あなたが乳がんである確率はどれぐらいでしょうか。
\end{itemize}
これは先ほどの公式に当てはめて考えると，
\[\frac{0.9\times0.01}{0.9\times0.01+0.08\times0.99}=0.102 \]

となり，10\%ちょっとだということになります\footnote{おっと，計算についていけないぞ！と思った人。大丈夫です，いい方法があります。確率の計算は小数点を含みますし相対的な度数なのでわかりにくいですから，絶対的な度数に戻して考えればいいのです。例えば10000人の世界だとして考えると，乳がんの人は1\%=100人です。乳がんでない人は9900人です。陽性になる人は乳がんで陽性反応が出るひと($100\times0.9=90$)と乳がんじゃないのに陽性反応が出る人($9900\times0.08=792$)の和で，882人です。陽性反応が出た人の中で本当に乳がんであるひとは，$90/882=0.102$となります。}。ずいぶんと印象が変わりますよね。PCR検査を複数回行うのは，これと同じような背景があるからです。

話を本題に戻しましょう。我々は心理学の研究の話をしているはずなのですから。

大事なのは，不確かな状況においては我々は確率的にしか判断をできず，ベイズの定理では条件の向きをひっくり返して計算できる，という点です。原因が結果を引き起こす確率が知りたい時に，結果を手にすることで，原因がどれほどあり得たかを考えることができる。それがベイズの定理のポイントなのです。

ベイズの定理の左辺には$P(A|B)$があり，右辺は分子に$P(B|A)$を含んでいます。$P(A)$(分子の2つ目の項目)や$P(B)$(分母)のことは一旦忘れてくれてもいいですし，
\[ P(A|B) = P(B|A) \times \frac{P(A)}{P(B)}\]
と考えて，$P(A|B)$が$P(B|A)$の関数である，と読み取ってくれても構いません。
ここでAを「心のモデル」，Bを「データ」と考えてみてください。

私たちは心の仕組みを考える時に，モデルをたてます。「あの人が$\diamond\diamond\diamond\diamond$な行動をとったのは，心の状態が$\dagger\dagger\dagger\dagger$だったからではないか」というのは心というモデルを立てて考えるということなのでした。心の状態が$\dagger\dagger\dagger\dagger$であれば，人は$\diamond\diamond\diamond\diamond$な行動をとるであろう，というのが研究者の仮説，モデルなのです。つまり，$P(\mbox{データ}|\mbox{モデル})$が心理学でやっていることなのです。

ところが実際には，手に入れることができるのはデータだけです。心がどうなっているかはわかりませんが，人が$\diamond\diamond\diamond\diamond$な行動をとった，という事実だけがそこにあるのです。この事実を使って，我々のモデルがどれぐらい正しいのか。あるいはモデルのパラメータはどうなっているのか，を検証したいわけです。つまり，研究の営みとは$P(\mbox{モデル}|\mbox{データ})$であり，データを所与(条件)としたときの，モデルの確率を考えることなのです。データをとってモデルを検証するという営みが成り立つ根拠を，ベイズの定理が示してくれていると言えるかもしれません\footnote{注意が必要なのは，モデルの正しい確率が検証されるわけではないところです。モデルを仮定すると，そのモデルから考えられる結果の確率がこれぐらいになる，という推論をしているのであり，モデルがそもそも前提されていますから，モデルが間違っていれば(例えばコイントスという現象に正規分布を当てはめて考えるとか)，その推論は全然正しいものではありません。}。

ベイズの定理を使ってパラメータを推論する時に，この定理が根本にあるということの意味が伝わりましたでしょうか。

